\section{Chapter~\ref{sec:eetypes}. Neurons as devices}

The operating modes of single cells are measured directly, with current injected through an electrode and voltage recording; the mathematics of the integrator and of the division into classes is classical theory, while the idea that the brain uses mode reconfiguration for computation remains a hypothesis.

{\footnotesize
\setlength{\tabcolsep}{3pt}\renewcommand{\arraystretch}{1.15}
\begin{longtable}{>{\raggedright}p{0.5cm}>{\raggedright}p{4.0cm}>{\raggedright}p{5.6cm}>{\raggedright}p{2.3cm}>{\raggedright\arraybackslash}p{1.7cm}}
\toprule
No. & Claim & Experiment and what was measured & Source & Status \\
\midrule\endfirsthead
\toprule
No. & Claim & Experiment and what was measured & Source & Status \\
\midrule\endhead
1 & Resonator: a slow current opposing voltage changes makes the neuron a band-pass filter peaking at a few to tens of hertz (Section~\ref{sec:resonator}) & In slices, sinusoidal current of smoothly rising frequency was injected into \nt{GLUT} neurons and the impedance $|Z(f)|$ was measured: a peak at $2$--$5$~Hz at $33^\circ$C (about $7$~Hz at $38^\circ$C); the resonance is created by slow potassium and cation currents and amplified by a persistent sodium current. A review describes the same technique for many cell classes & \cite{HuVervaekeStorm2002,HutcheonYarom2000} & measured \\
2 & The slow current behaves like an inductance and, together with the membrane capacitance, forms a tank circuit & The subthreshold response of an axon to current was measured and described by linearized conductance equations; the equivalent circuit contains a ``phenomenological'' inductance whose value depends on voltage & \cite{Mauro1970} & theory \\
3 & The time constant of a group with feedback is $\tau_{\text{eff}}=\tau/(1-w)$~\eqref{eq:perfectint}; with $\tau=0.1$~s and $w=0.995$ it is $20$~s & Derived in the chapter from the linear group equation; proposed in a theoretical paper as the mechanism for holding eye position & derived in chapter; \cite{Seung1996} & theory \\
4 & The eye-position integrator holds its input by the network, not by a property of a single cell & Intracellular recording in fish in the nucleus that holds eye position: after a fast eye turn the neuron's rate changes in a step and holds; a brief current into one cell produces only a transient shift, while the voltage steps persist below threshold as well, sustained by synaptic inputs & \cite{Aksay2001} & measured \\
5 & A leak-free integrator needs constant retuning by learning; when mistuned, it either forgets or runs into saturation & Fish were trained with visual feedback proportional to $\pm$ eye position: holding became unstable or leaky, and the neurons' time constant fell by more than an order of magnitude, to $<1$~s; ordinary visual feedback gradually restored stability & \cite{Major2004} & measured \\
6 & Bistable neuron: a brief input switches on a lasting plateau, inhibition switches it off; the property is enabled by serotonin (Section~\ref{sec:bistable}) & Intracellular recording from the \nt{ACET} neurons that control muscles in the cat: brief synaptic excitation puts the neuron into lasting activity, brief inhibition resets it; in the spinal cat the state appears after a serotonin precursor is given & \cite{Hounsgaard1988} & measured \\
7 & Two classes: integrators (rate rises from zero) and resonators (a finite rate right at threshold) & The classes are derived from bifurcation theory. Experiment: in cortical slices, \nt{GLUT} neurons start firing at arbitrarily low rates (type~1), fast \nt{GABA} neurons jump in at a finite rate (type~2) & \cite{Izhikevich2007,Tateno2004} & measured \\
8 & One neuron is an integrator or a coincidence detector: inhibition shortens the summation window & In hippocampal slices, the window within which excitatory inputs sum to threshold is shorter than $2$~ms; it is shortened by feedforward inhibition arriving right after the excitation; in dendrites, where inhibition is weaker, the window is wider & \cite{Pouille2001} & measured \\
9 & A delay line made of axons (Table~\ref{tab:eetypes}) & In the barn owl, delays were measured in the axons leading to coincidence detectors: different axon lengths give different delays, and the detectors are tuned to the difference in sound arrival times & \cite{Carr1990} & measured \\
10 & A pacemaker in wakefulness and a silent cell in sleep & \nt{SERO} neurons fire almost regularly during the day, slow down in slow-wave sleep, and are nearly silent in REM sleep (more in Chapter~\ref{sec:sleep}) & \cite{JacobsAzmitia1992,McGintyHarper1976} & measured \\
11 & The device is set by ion channels and by the broadcast channels that tune them (Proposition~\ref{prop:eetypes}) & Shown in small networks: modulatory substances change the intrinsic properties of neurons and the strength of synapses, so the same circuit produces different rhythms & \cite{Marder2012} & measured \\
12 & The brain uses mode reconfiguration for computation (Proposition~\ref{prop:eetypes}) & There is no direct experiment in which reconfiguring a cell's mode was needed to solve a task; there are only experiments in which modulators change a circuit's output & \cite{Marder2012} & hypothesis \\
\bottomrule
\end{longtable}}
